\section{Holomorphic sectional curvature}
\label{sec:integrated-hermitian-curvature}

The curvature conditions above test a tangent direction and a fiber vector separately. For a metric on the tangent bundle, holomorphic sectional curvature tests the same vector in both places. Averaging these diagonal evaluations controls a symmetrized curvature trace. A second trace occurs when one differentiates the trace of the metric itself; their difference contains torsion and an exact form. With a K\"ahler metric for the contraction, integration removes the exact form and the torsion terms cancel. This gives an obstruction even when the metric whose curvature is tested is not closed.

In this section, $g$ is a K\"ahler metric on $X$ and $h$ is any smooth Hermitian metric on $T^{1,0}X$. In local holomorphic coordinates, their forms are $\omega_g=\mathrm i\sum_{k,l}g_{k\bar l}\,\mathrm dz^k\wedge\mathrm d\bar z^l$ and $\omega_h=\mathrm i\sum_{k,l}h_{k\bar l}\,\mathrm dz^k\wedge\mathrm d\bar z^l$. As in the rest of the chapter, $n\geq2$. The unscaled Chern curvature convention of \cref{subsec:projective-complement-curvature-index} gives the coefficients
\[
 R^h_{i\bar j k\bar l}
 =-\partial_i\partial_{\bar j}h_{k\bar l}
   +h^{p\bar q}(\partial_i h_{k\bar q})
                       (\partial_{\bar j}h_{p\bar l}).
\]
The \emph{Chern holomorphic sectional curvature} is
\begin{equation}\label{eq:chern-holomorphic-sectional-curvature}\index{Holomorphic sectional curvature!Chern connection}
 H_h(x,[v])=\frac{R^h(v,\bar v,v,\bar v)}{|v|_h^4},
 \qquad 0\ne v\in T_x^{1,0}X.
\end{equation}
It is \emph{quasi-negative} if it is nonpositive everywhere and strictly negative in every direction at some point. Put
\[
 u=\operatorname{tr}_g h,\qquad
 q=\operatorname{tr}\bigl((g^{-1}h)^2\bigr).
\]
In a $g$-unitary frame, define the real contractions
\[
 Q_g(h)=\sum_{i,k}R^h_{i\bar i k\bar k},\qquad
 B_g(h)=\sum_{i,k}R^h_{i\bar k k\bar i},\qquad
 \mathcal R_g(h)=\tfrac12\bigl(Q_g(h)+B_g(h)\bigr).
\]
These expressions are invariant under a change of unitary frame. For a general $h$, the two traces $Q_g(h)$ and $B_g(h)$ can differ.

\Needspace{10\baselineskip}
\begin{lemma}\label{lem:hermitian-sectional-average}\index{Chern curvature!symmetrized trace}
At a point where $H_h\leq-\kappa$ for some $\kappa\geq0$, one has
\[
 \mathcal R_g(h)\leq-\frac{\kappa}{2}(u^2+q).
\]
In particular, nonpositive holomorphic sectional curvature gives $\mathcal R_g(h)\leq0$.
\end{lemma}
\begin{proof}
In a $g$-unitary frame, let $Z_1,\ldots,Z_n$ be independent standard complex Gaussian variables, with density $\pi^{-1}\mathrm e^{-|z|^2}$ on $\mathbb C$. Their second and fourth moments are $\mathbb E|Z_i|^2=1$ and $\mathbb E|Z_i|^4=2$. Independence and invariance under coordinate phases give
\[
 \mathbb E[Z_i\overline{Z_j}Z_k\overline{Z_l}]
      =\delta_{ij}\delta_{kl}+\delta_{il}\delta_{kj}.
\]
Thus $\mathbb E R^h(Z,\bar Z,Z,\bar Z)=Q_g(h)+B_g(h)$ and
$\mathbb E|Z|_h^4=u^2+q$. Taking expectations in the homogeneous inequality
$R^h(Z,\bar Z,Z,\bar Z)\leq-\kappa|Z|_h^4$ proves the result.
\end{proof}

Let $A=\nabla^h-\nabla^g$ be the difference of the Chern connections, viewed as a tensor with two holomorphic covariant indices and a tangent-bundle value. Write
\[
 S(v,w)=\tfrac12\bigl(A(v,w)+A(w,v)\bigr),\qquad
 T(v,w)=A(v,w)-A(w,v).
\]
Since $g$ is K\"ahler, $T$ is the Chern torsion of $h$. Norms such as
$|S|_{g,h}^2=\sum_{i,k}|S(e_i,e_k)|_h^2$, where $(e_i)$ is a $g$-unitary frame, use $g$ in the covariant indices and $h$ in the value. Let $r_{i\bar j}$ be the coefficients of the Ricci form
$\operatorname{Ric}(g)=-\mathrm i\partial\bar\partial\log\det(g_{i\bar j})$, and in a $g$-unitary frame set
\[
 I_g(h)=\sum_{i,k}r_{i\bar k}h_{k\bar i},\qquad
 \mathrm dV_g=\frac{\omega_g^n}{n!},\qquad
 P_g f=\operatorname{tr}_{\omega_g}(\mathrm i\partial\bar\partial f).
\]

\Needspace{10\baselineskip}
\begin{proposition}\label{prop:integrated-hermitian-chern-lu}\index{Chern curvature!integrated trace identity}
The two metrics satisfy
\begin{equation}\label{eq:integrated-hermitian-chern-lu}
 \int_X\bigl(|S|_{g,h}^2+I_g(h)-\mathcal R_g(h)\bigr)\,\mathrm dV_g=0.
\end{equation}
\end{proposition}
\begin{proof}
We first compute the exact term that distinguishes the two curvature traces. Define
\begin{align*}
 D_{i\bar j k\bar l}
 ={}&\partial_i\partial_{\bar j}h_{k\bar l}
     +\partial_k\partial_{\bar l}h_{i\bar j}\\
   &-\partial_i\partial_{\bar l}h_{k\bar j}
     -\partial_k\partial_{\bar j}h_{i\bar l}.
\end{align*}
This is a tensor: putting holomorphic differentials first gives
\[
 \mathrm i\partial\bar\partial\omega_h
 =\sum_{i<k,\,j<l}D_{i\bar j k\bar l}
          \,\mathrm dz^i\wedge\mathrm dz^k
             \wedge\mathrm d\bar z^j\wedge\mathrm d\bar z^l.
\]
For $D_g=\sum_{i,k}D_{i\bar i k\bar k}$ in a $g$-unitary frame, expansion of the wedge product gives
\[
 D_g\,\mathrm dV_g
 =2\mathrm i\partial\bar\partial\omega_h
                    \wedge\frac{\omega_g^{n-2}}{(n-2)!}.
\]
Indeed, each pair $i<k$ contributes $D_{i\bar i k\bar k}\mathrm dV_g$ to the wedge product and occurs twice in the contraction. Closedness of $\omega_g$ and Stokes' theorem therefore give $\int_XD_g\,\mathrm dV_g=0$.

The four-term antisymmetrization of the curvature formula is
\begin{align*}
 &R^h_{i\bar j k\bar l}+R^h_{k\bar l i\bar j}
  -R^h_{i\bar l k\bar j}-R^h_{k\bar j i\bar l}\\
 &\hspace{3em}=h_{p\bar q}T^p_{ik}\overline{T^q_{jl}}
                         -D_{i\bar j k\bar l}.
\end{align*}
The second derivatives give $-D_{i\bar jk\bar l}$. The quadratic term of $R^h_{i\bar jk\bar l}$ is $h_{p\bar q}\Gamma^p_{ik}\overline{\Gamma^q_{jl}}$, where $\Gamma^p_{ik}=h^{p\bar q}\partial_ih_{k\bar q}$ are the Chern connection coefficients of $h$. Since $g$ is K\"ahler, $\Gamma^p_{ik}-\Gamma^p_{ki}=T^p_{ik}$, so the four quadratic terms combine to $h_{p\bar q}T^p_{ik}\overline{T^q_{jl}}$. Contracting gives
$2(Q_g(h)-B_g(h))=|T|_{g,h}^2-D_g$, hence
\[
 Q_g(h)=\mathcal R_g(h)+\tfrac14|T|_{g,h}^2-\tfrac14D_g.
\]

Now choose holomorphic coordinates that are normal for $g$ at a given point, so that $g$ is the identity and its first derivatives vanish there; then $A$ has the coefficients $\Gamma^p_{ik}$ at this point. When $u=g^{k\bar l}h_{k\bar l}$ is differentiated twice, the second derivatives of the inverse metric contribute $I_g(h)$, by the K\"ahler curvature symmetries of $g$, while the second derivatives of $h$ contribute $-Q_g(h)+|A|_{g,h}^2$. Thus
\begin{align*}
 P_g u&=|A|_{g,h}^2+I_g(h)-Q_g(h)\\
      &=|S|_{g,h}^2+I_g(h)-\mathcal R_g(h)+\tfrac14D_g,
\end{align*}
where symmetric and skew tensors are orthogonal and
$|A|_{g,h}^2=|S|_{g,h}^2+|T|_{g,h}^2/4$. Integrating removes both $P_g u$ and $D_g$, proving \eqref{eq:integrated-hermitian-chern-lu}.
\end{proof}

\Needspace{10\baselineskip}
\begin{corollary}\label{cor:nonnegative-ricci-sectional-obstruction}\index{Holomorphic sectional curvature!Ricci obstruction}
Suppose that $\operatorname{Ric}(g)\geq0$. Every Hermitian tangent metric $h$ with $H_h\leq0$ satisfies $\operatorname{Ric}(g)=0$ and $H_h=0$ in every direction. In particular, $X$ admits no Hermitian tangent metric with quasi-negative Chern holomorphic sectional curvature.
\end{corollary}
\begin{proof}
The Ricci assumption gives $I_g(h)\geq0$, and \cref{lem:hermitian-sectional-average} gives $-\mathcal R_g(h)\geq0$. All three continuous terms in \eqref{eq:integrated-hermitian-chern-lu} are nonnegative, so each vanishes everywhere. A positive definite $h$ pairs strictly positively with every nonzero semipositive Ricci matrix; hence $I_g(h)=0$ forces $\operatorname{Ric}(g)=0$.

At each point, the Gaussian expectation in the averaging proof is now zero. Its integrand $R^h(Z,\bar Z,Z,\bar Z)$ is nonpositive. If it were negative in one direction, continuity and the strictly positive Gaussian density would make its expectation negative. It is therefore zero in every direction, which gives $H_h=0$ and excludes quasi-negativity.
\end{proof}

For the Ricci-flat surfaces of \cref{sec:ricci-flat-surface-rc}, this obstruction concerns all Hermitian tangent metrics, including nonclosed ones. It coexists with the RC-positive metric of \cref{thm:ricci-flat-surface-rc-positive}: that condition allows the positive tangent direction to depend on the fiber vector, while holomorphic sectional curvature evaluates the same direction twice.

\subsection{Rigidity on complex tori}
\label{subsec:torus-sectional-rigidity}

For a flat reference metric, the vanishing symmetric connection term has a stronger consequence. Periodic integration by parts forces the second metric to be translation invariant, and hence to have zero full Chern curvature.

\Needspace{12\baselineskip}
\begin{proposition}\label{prop:torus-sectional-rigidity}\index{Complex torus!nonpositive holomorphic sectional curvature}\index{Holomorphic sectional curvature!rigidity on a torus}
Let $Y=\mathbb C^m/\Lambda$ be a complex torus, with $m\geq1$. A smooth Hermitian tangent metric $h$ on $Y$ has nonpositive Chern holomorphic sectional curvature if and only if it is translation invariant. Such a metric is K\"ahler and has zero Chern curvature.
\end{proposition}
\begin{proof}
Suppose first that $m\geq2$, and take the flat K\"ahler metric $g$ induced by a constant Hermitian inner product on $\mathbb C^m$. Choose linear coordinates in which this inner product is the identity. The coordinate vector fields descend to $Y$, even though the coordinates themselves need not. By \cref{cor:nonnegative-ricci-sectional-obstruction}, $H_h=0$; \eqref{eq:integrated-hermitian-chern-lu} also gives $S=0$. Since $g$ has zero connection coefficients, multiplying $S=0$ by $h$ gives
\[
 \partial_i h_{k\bar l}+\partial_k h_{i\bar l}=0
 \qquad\text{for all }i,k,l.
\]
Fix $l$ and put $a_k=h_{k\bar l}$. With respect to the flat volume, integration by parts and commutation of constant coordinate derivatives give
\[
 \sum_{i,k}\int_Y\partial_i a_k\,
          \overline{\partial_k a_i}\,\mathrm dV_g
 =\int_Y\left|\sum_k\partial_{\bar k}a_k\right|^2\mathrm dV_g.
\]
Consequently,
\[
 \begin{aligned}
 0&=\sum_{i,k}\int_Y|\partial_i a_k+\partial_k a_i|^2\,\mathrm dV_g\\
  &=2\sum_{i,k}\int_Y|\partial_i a_k|^2\,\mathrm dV_g
    +2\int_Y\left|\sum_k\partial_{\bar k}a_k\right|^2\mathrm dV_g.
 \end{aligned}
\]
Thus every $\partial_i h_{k\bar l}$ vanishes. Hermitian symmetry gives the vanishing of the conjugate derivatives as well, so all coefficients of $h$ are constant.

For $m=1$, write the lifted metric as $h(z)|\mathrm dz|^2$, with $h$ positive and $\Lambda$-periodic. Its curvature coefficient is $-h\,\partial_z\partial_{\bar z}\log h$. Nonpositive sectional curvature therefore gives $\partial_z\partial_{\bar z}\log h\geq0$. The complex Laplacian has zero integral over $Y$, so it vanishes. Integration by parts then shows that $\log h$ has zero gradient, and $h$ is constant.

Conversely, a constant positive Hermitian matrix defines a closed form and has zero connection and curvature coefficients. This proves all the assertions.
\end{proof}

\Needspace{9\baselineskip}
\begin{corollary}\label{cor:torus-cover-sectional-rigidity}\index{Holomorphic sectional curvature!unramified torus covers}
Let $Y$ be a connected compact complex manifold admitting a finite unramified holomorphic covering $\pi:Z\to Y$ by a complex torus. Every smooth Hermitian tangent metric on $Y$ with nonpositive Chern holomorphic sectional curvature is K\"ahler and Chern flat.
\end{corollary}
\begin{proof}
The covering is locally biholomorphic, so the pullback metric has the same sectional-curvature sign. By \cref{prop:torus-sectional-rigidity}, its associated form is closed and its Chern curvature is zero. Both conclusions descend through the covering.
\end{proof}

